\section{Related Work}

Work on edge and fog data reduction consistently treats filtering, aggregation, and adaptation as lightweight mechanisms for reducing communication pressure without relocating the full analytics stack into the network. Surveys by Pioli et al.\ and Sadri et al.\ emphasize both the practical appeal of such mechanisms and the need for evaluations that report trade-offs explicitly rather than assuming universal gains \cite{pioli2022_edge_data_reduction_mapping,sadri2022_fog_data_reduction_survey}. Agrasandhani follows this line of work at the gateway layer, but evaluates the policies against dashboard-facing outcomes: downstream frame cadence, display latency, and visible age under last-hop impairment.

MQTT-specific work motivates the treatment of delivery semantics, duplicates, and freshness. The MQTT 3.1.1 specification distinguishes QoS~0 and QoS~1 delivery guarantees, and Kovalchuk et al.\ note that at-least-once delivery can require application-side handling of duplicate messages \cite{mqtt311,kovalchuk2018_mqtt_architecture_delay}. Palmese et al.\ argue for adaptation under changing network conditions, while Yates et al.\ frame freshness as an explicit observer-side property rather than a by-product of throughput \cite{palmese2022_adaptive_qos_mqttsn,yates2021_age_of_information_survey}. These ideas align with Agrasandhani's batch-window control, compaction, and age-aware dashboard semantics, even though the measured adaptive result here remains negative under the tested thresholds.

Methodologically, the paper is closer to controlled impairment studies than to open-ended deployment measurement. Jurgelionis et al.\ and Garcia and Hurtig argue that deterministic impairment is valuable because protocol and application behavior can then be compared under repeatable degradations \cite{jurgelionis2011_netem_empirical,garcia2016_kaunetem_netdev}. Agrasandhani adopts that methodological stance by using a deterministic application-layer impairment proxy on the gateway-to-dashboard path and by reporting proxy-side outputs as the canonical source of truth for dashboard-facing stability.
